\documentclass[11pt]{article}

% ====================================================
% PACKAGES
% ====================================================
\usepackage[margin=1in, headheight=13.6pt]{geometry}
\usepackage[T1]{fontenc}
\usepackage{lmodern}
\usepackage{microtype}
\usepackage{setspace}
\usepackage{tcolorbox}
\usepackage{amsmath, amssymb, mathtools}
\usepackage{siunitx}

\usepackage{graphicx}
\usepackage{float}
\usepackage{subcaption}

\usepackage{enumitem}
\usepackage{booktabs}
\usepackage{longtable}
\usepackage{array}
\usepackage{multirow}

\usepackage{xcolor}
\usepackage{hyperref}

\usepackage{fancyhdr}
\usepackage{lastpage}

\usepackage{listings}
\usepackage{indentfirst}
\usepackage[justification=centering]{caption}

% Python source style for \CodeListing
\lstdefinestyle{pystyle}{
    language=Python,
    basicstyle=\ttfamily\footnotesize,
    keywordstyle=\color{blue!60!black}\bfseries,
    stringstyle=\color{orange!70!black},
    commentstyle=\color{gray!60!black}\itshape,
    numbers=left,
    numberstyle=\tiny\color{gray!70},
    numbersep=8pt,
    breaklines=true,
    showstringspaces=false,
    frame=single,
    rulecolor=\color{gray!40},
    tabsize=4,
    columns=fullflexible,
    captionpos=b,
    literate={°}{{\textdegree}}1,
}
\lstset{style=pystyle}

\pagestyle{fancy}
\fancyhf{}
\lhead{\Assignment\ |\ \course}
\rhead{\Name\ |\ \SubDate}
\cfoot{Page \thepage\ of~\pageref{LastPage}}

\hypersetup{
    colorlinks,
    linkcolor={red!50!black},
    citecolor={blue!50!black},
    urlcolor={blue!80!black}
}

\newcommand{\Name}{Arael Anaya}
\newcommand{\Assignment}{Initial Project Summary}
\newcommand{\course}{Advanced Dynamics}
\newcommand{\SubDate}{7/20/26}

% Boxed final-answer environment
\newtcolorbox{result}{
    colback=blue!3!white,
    colframe=blue!40!black,
    boxrule=0.6pt,
    arc=2pt,
    left=6pt, right=6pt, top=4pt, bottom=4pt
}

% Code listing: sources the .py file directly.
\newcommand{\CodeListing}[3]{%
\IfFileExists{#1}%
    {\lstinputlisting[style=pystyle,caption={#2},label={#3}]{#1}}%
    {\begin{quote}\itshape Source file \texttt{#1} not found yet --- this
        listing will populate automatically once the file is written.\end{quote}}}

\begin{document}

% ====================================================
% HEADER BLOCK (matches original submission)
% ====================================================
\begin{flushright}
    \Name \\
    \course \\
    \SubDate
\end{flushright}

\vspace{0.5em}

% ====================================================
% PROJECT SUMMARY
% ====================================================
For my final project I plan to study the spring-loaded inverted pendulum
(SLIP) model of a hopping/running robot, which is the base for essentially
all legged robot models. Despite its initial simplicity (a point mass
bouncing on a massless spring leg), the SLIP model captures the
center-of-mass dynamics of running in both animals and machines shockingly
well. That makes it a great subject for this project: it is simple enough to
analyze carefully with the tools from this course, but rich enough to
produce very interesting behavior. It also builds naturally from the
elastic pendulum system we have analyzed in this classroom, since the
stance phase is essentially an inverted spring pendulum.

What makes the SLIP model especially interesting from a mechanics
standpoint is that it is a hybrid dynamical system. During the stance phase
the body behaves like a spring-loaded pendulum, while in the flight phase
it is just a projectile. I will derive the equations of motion for each
phase, using the Lagrangian formulation in polar coordinates for the stance
phase, and work out the switching conditions at touchdown and liftoff that
connect the two phases.

Rather than just analyzing the stability at a single equilibrium point, I
will study the stability of the hopping gait itself. From quick research,
the natural tool for this is an apex return map (a Poincar\'e map in which
a periodic gait appears as a fixed point whose stability can be evaluated
numerically). I will implement this model in Python, using event-based
numerical integration, and run a parameter study on quantities like spring
stiffness and touchdown angle to map out which combinations produce stable
hopping. There is substantive literature on SLIP dynamics, so there is
plenty to draw from for the paper and video presentation.

\end{document}
